Water quality indices condense several physical and chemical measurements into a single
interpretable score through a fixed formula. This creates a problem when such an index
becomes a supervised learning target: a model trained on the measurements that define it
recovers the formula rather than learning anything about water quality. We demonstrate
this on 2{,}818{,}006 harmonised surface-water records spanning five countries and eight
decades, where gradient boosting reaches $R^2 = 0.9997$ against the CCME Water Quality
Index -- a figure that reflects function approximation, not prediction. We therefore
reframe the task as estimating the index when measurements are unavailable, the case in
which the formula cannot be evaluated and an estimate is genuinely required. Across four
availability regimes, performance falls to $R^2 = 0.834$ when one measurement is withheld
and $0.573$ when four are, with mean absolute error rising from 0.100 to 8.596 index
points. XGBoost leads Random Forest at every level on identical site-grouped test
data, though on a larger training budget.
We further identify an upstream median-imputation artefact affecting up to 31.6\% of
individual columns, concentrated almost entirely in one country and confounded with the
target, and show by ablation that it does not warrant inclusion as a feature.
